\documentclass{article}
\usepackage[margin=1.15in]{geometry}
\usepackage{xcolor}
\usepackage{parskip}
\usepackage{fancyhdr}
\usepackage{array}
\usepackage{booktabs}
\usepackage{enumitem}
\usepackage{amsmath}
\usepackage{tabularx}
\usepackage{listings}

\pagestyle{fancy}
\fancyhf{}
\cfoot{\thepage}
\rhead{Lecture 16}
\lhead{MA 214: Applied Statistics (Spring 2026)}
\renewcommand{\headrulewidth}{.4mm}

\newcommand{\blankline}[1]{\underline{\hspace{#1}}}

\lstset{
  basicstyle=\ttfamily\small,
  columns=fullflexible,
  frame=single,
  framerule=0.4pt,
  breaklines=true,
  showstringspaces=false
}

\begin{document}

\noindent\textbf{Name:}\blankline{2.25in}\hfill\textbf{BUID:}\blankline{1.55in}

\medskip

\noindent\textbf{Name:}\blankline{2.25in}\hfill\textbf{BUID:}\blankline{1.55in}\newline

\begin{center}
 \Large In-Class Worksheet: Logistic Regression Inference
\end{center}

\subsection*{Scenario}
The \emph{Cleveland heart-disease} study (Detrano et al., 1989) records $n = 208$ patients. The predictors are \texttt{age}, \texttt{sex}, \texttt{cholesterol}, \texttt{resting\_bp}, \texttt{max\_heart\_rate}, \texttt{exercise\_angina} and \texttt{fasting\_blood\_sugar}; the outcome is a binary \texttt{heart\_disease} diagnosis. You will fit a logistic model, do inference on the coefficients, check the conditions, and compare models by cross-validation. Data (\texttt{heart\_disease.csv}) is on Blackboard.

\subsection*{Part 1: Explore}

\begin{enumerate}[label={\bfseries (\Alph*)},itemsep=.7em]
\item Load the data and record: \; number of patients \blankline{0.9in}; \; proportion with heart disease \blankline{0.9in}; \; mean age \blankline{0.9in}. \\[0.5em]
\item Plot \texttt{heart\_disease} (jittered $0/1$) against \texttt{age} with a fitted logistic curve:
\begin{lstlisting}[language=R]
geom_smooth(method = "glm", method.args = list(family = "binomial"))
\end{lstlisting}
Does the probability of heart disease rise or fall with age? \\[3em]
\end{enumerate}

\subsection*{Part 2: Fit and read the output}

Fit the model:
\begin{lstlisting}[language=R]
glm(heart_disease ~ age + resting_bp + cholesterol + max_heart_rate,
    family = binomial)
\end{lstlisting}

\begin{enumerate}[label={\bfseries (\Alph*)},itemsep=.7em]
\item From \texttt{tidy()}, record the coefficient table:
\vspace{0.2em}
\begin{center}
\renewcommand{\arraystretch}{1.3}
\begin{tabular}{lcccc}
\toprule
\textbf{term} & \textbf{estimate} & \textbf{std.error} & \textbf{$z$} & \textbf{p.value} \\
\midrule
(Intercept)    & & & & \\
age            & & & & \\
resting\_bp    & & & & \\
cholesterol    & & & & \\
max\_heart\_rate & & & & \\
\bottomrule
\end{tabular}
\end{center}
\item Which predictors show strong evidence of association ($p < 0.05$)? Interpret the \texttt{age} coefficient on the \emph{log-odds} scale. \\[3.5em]
\item Run \texttt{tidy(conf.int = TRUE, exponentiate = TRUE)}. Interpret the \texttt{resting\_bp} odds ratio \emph{in context} --- and per a meaningful step (say, $10$ mmHg), not per single unit. \\[4em]
\end{enumerate}

\subsection*{Part 3: Confidence intervals and diagnostics}

\begin{enumerate}[label={\bfseries (\Alph*)},itemsep=.7em]
\item Does the $95\%$ CI for the \texttt{age} odds ratio include $1$? What does that say about the evidence for age? \\[3em]
\item Does the CI for \texttt{max\_heart\_rate} include $1$? What does that say? \\[3em]
\item Make the predicted-probability histogram colored by outcome. Do the two groups separate? \\[3em]
\item Make a binned-residual plot. Are any bins outside the $\pm 2\,\mathrm{SE}$ band? What would that indicate? \\[3.5em]
\end{enumerate}

\subsection*{Part 4: Prediction and cross-validation}

\begin{enumerate}[label={\bfseries (\Alph*)},itemsep=.7em]
\item Classify at $c = 0.5$ and fill in the confusion matrix; then the misclassification rate.
\vspace{0.2em}
\begin{center}
\begin{tabular}{lcc}
\toprule
& \textbf{Predicted 1} & \textbf{Predicted 0} \\
\midrule
\textbf{Actual 1} & TP = \blankline{1.4cm} & FN = \blankline{1.4cm} \\[0.3em]
\textbf{Actual 0} & FP = \blankline{1.4cm} & TN = \blankline{1.4cm} \\
\bottomrule
\end{tabular}
\end{center}
Misclassification rate $= \blankline{2cm}$
\item Compute the Brier score \blankline{1.3in} and log-loss \blankline{1.3in} on the full data. In one sentence, how do the two penalise a confident wrong prediction differently? \\[3em]
\item \textbf{Cross-validate.} Run $5$-fold CV (repeat a few times and average) on the four models below, scoring by misclassification, Brier, and log-loss.
\vspace{0.2em}

\begin{center}\small
\renewcommand{\arraystretch}{1.5}
\begin{tabular}{@{}cl ccc@{}}
\toprule
\textbf{\#} & \textbf{Model (predictors)} & \textbf{Misclass} & \textbf{Brier} & \textbf{LogLoss} \\
\midrule
1 & age + resting\_bp + cholesterol + max\_hr & \blankline{0.8cm} & \blankline{0.8cm} & \blankline{0.8cm} \\
2 & Model 1 $+$ fasting\_blood\_sugar & \blankline{0.8cm} & \blankline{0.8cm} & \blankline{0.8cm} \\
3 & Model 1 $+$ exercise\_angina & \blankline{0.8cm} & \blankline{0.8cm} & \blankline{0.8cm} \\
4 & Model 1 $+$ both & \blankline{0.8cm} & \blankline{0.8cm} & \blankline{0.8cm} \\
\bottomrule
\end{tabular}
\end{center}
\vspace{0.4em}

Do the three losses agree on the best model? \blankline{1.5in} \; Which would you choose, and why? \\[3.5em]
\item \textbf{Threshold.} Sweep the threshold from $0.1$ to $0.9$. \\[0.3em] Where is misclassification lowest? \blankline{1in} \\[0.3em] For a heart-disease \emph{screen}, would a doctor want a higher or lower one? Why? \\[3.5em]
\end{enumerate}

\newpage

\subsection*{Part 5: Wrap-up}

\begin{enumerate}[label={\bfseries (\Alph*)},itemsep=.8em]
\item Name two ways inference in logistic regression differs from inference in linear regression. \\[4em]
\item Why might cross-validation with the Brier score or log-loss rank the models differently from AIC or the misclassification rate? \\[4em]
\end{enumerate}

\subsection*{Compare with the groups around you}

\noindent\fbox{\parbox{\dimexpr\textwidth-2\fboxsep-2\fboxrule\relax}{%
Compare your final model with at least two neighboring pairs. Did everyone read the cholesterol odds ratio the same way --- per unit, or per meaningful step? For a screening test, would you set the threshold above or below $0.5$? Having heard them, would you change your model or your threshold?
\\[5em]
}}

\end{document}
